\documentclass[12pt, oneside, paper=A4, DIV=15, BCOR=0mm, abstract=true, headings=small]{scrartcl}
\usepackage[english]{babel}
\usepackage[utf8]{inputenc}
\usepackage{lmodern}

\usepackage[T1]{fontenc}

\usepackage[footsepline]{scrlayer-scrpage}

\usepackage[auto]{microtype}
\clubpenalty = 10000
\widowpenalty = 10000
\displaywidowpenalty = 10000

\usepackage{graphicx}
\usepackage{multirow, multicol, booktabs}
\usepackage{threeparttable}
\usepackage{longtable}
\usepackage{rotating}
\usepackage{ltablex}
\usepackage{subfig}
\captionsetup[subtable]{position=top}
\usepackage{pdfpages}
\usepackage{listings}
\usepackage{comment}
\usepackage{appendix}

\usepackage{url}
\urlstyle{same}

\usepackage{footnote}
\makesavenoteenv{tabular} 

\usepackage{todonotes}
\usepackage{blindtext}

\usepackage[
backend=biber,
style=iso-numeric,
citestyle=numeric-comp,
maxbibnames=2,
firstinits=true
]{biblatex}

\renewcommand*{\labelnamepunct}{\addcolon\addspace}

\bibliography{Bibliography/Biblio.bib}

\deffootnote{1.5em}{1em}{%
	\makebox[1.5em][l]{\thefootnotemark}%
}

\interfootnotelinepenalty=10000

\addtokomafont{caption}{\small}
\setkomafont{captionlabel}{\sffamily\bfseries}
\setkomafont{subject}{\Large\bfseries}
\setkomafont{author}{\normalfont}
\setkomafont{date}{\normalfont}
\setkomafont{publishers}{\normalfont}


\makeatletter

\renewcommand{\l@section}{\@dottedtocline{1}{0ex}{1.8em}}
\renewcommand{\l@subsection}{\@dottedtocline{2}{1.8em}{2.7em}}

\makeatother

\newenvironment{tabular10}{%
	\fontsize{10}{12}\selectfont\tabular
}{%
	\endtabular
}

\newenvironment{tabular7}{%
	\fontsize{7}{12}\selectfont\tabular
}{%
	\endtabular
}

\usepackage[
pdftitle={Abschlussbericht},
pdfsubject={},
pdfauthor={},
pdfkeywords={},  
hidelinks
]{hyperref}

% Formato nel testo:
% Lecture 0. Introduction

\renewcommand*{\sectionformat}{%
	Lecture~\thesection.\enskip
}

% Formato nell'indice:
% Lecture 0. Introduction ........ 2

\makeatletter

\newcommand*{\sectionentrynumberformat}[1]{%
	\begingroup
	% Inserisce il punto subito dopo il numero,
	% prima dello spazio elastico della numberline
	\renewcommand*{\numberline@numberpostfix}{.}%
	Lecture~#1%
	\endgroup
}

\makeatother

\DeclareTOCStyleEntry[
indent=0pt,
numwidth=6em,
entrynumberformat=\sectionentrynumberformat,
linefill=\TOCLineLeaderFill
]{tocline}{section}
